\begin{tikzpicture}[scale=1]
\tikzmath{
\R=1.5;
\f=0.5;
\fx=\f+1;
}
\scalebox{1.0}{

\coordinate (P) at (0.8*\R,0.7*\R);
\coordinate (PP) at (0.8*\R,-0.5*\R);
\coordinate (PX) at (0.8*\R*\fx,0);
\coordinate (PZ) at (-\f*0.8*\R,-0.5*\R);
 %axes:
 \draw[-latex, line width=1.5pt] (0,0) -- (1.75*\R,0) node[below]{$x$}; 
 \draw[-latex, line width=1.5pt] (0,0) -- (0,1.75*\R) node[left]{$y$}; 
 \draw[-latex, line width=1.5pt] (0,0) -- (-\f*1.75*\R,-\f*1.75*\R) node[left]{$z$};  
 
 %point
 \fill (P) circle(2pt);
 \node[right] at (P) {$P$};
 
 %dotted lines
 \draw[dashed,color=red,line width=1.2pt] (P) -- (PP) node[midway,right]{$y_P$};
 \draw[dashed,color=red,line width=1.2pt] (PP) -- (PX) node[midway,right]{$z_P$};
 \draw[dashed,color=red,line width=1.2pt] (PP) -- (PZ) node[midway,below]{$x_P$};
}
\end{tikzpicture}
\captionof*{figure}{\textit{A \textbf{right-handed} Cartesian coordinate system, where $\hat x \times \hat y = \hat z$. ($\hat z$ is out of the page). The point $P$ has coordinates $(x_P,y_P,z_P)$.}}